%=============================================================================
%  Chapitre 7 — Sprint 4 : Systèmes de Sécurité Gaz et DHT11
%=============================================================================
\chapter{Sprint 4 --- Détection Gaz, Capteur DHT11 et Protocoles d'Urgence}

\begin{tcolorbox}[sprintbox,title={Sprint 4 --- Systèmes de Sécurité}]
\textbf{Durée :} 2 semaines \quad
\textbf{Points de Story :} 36 \quad
\textbf{Objectif :} Implémenter la détection d'urgence gaz de bout en bout depuis le
capteur MQ6 jusqu'à l'actuation matérielle et l'alerte UI plein écran, plus la
surveillance de température et humidité DHT11.
\end{tcolorbox}

\section{Service de Fond Gas Monitor}

Le service \texttt{gasMonitor.js} est le cerveau sécuritaire de la plateforme. Il
s'abonne à tous les topics de capteurs gaz, évalue les lectures PPM par rapport au seuil
configuré, et pilote la machine à états d'urgence.

\begin{lstlisting}[style=codeblock,language=JavaScript,caption={Service Gas Monitor (backend/gasMonitor.js)}]
const initGasMonitor = (mqttClient, io) => {
  mqttClient.subscribe('sensors/+/gas');

  mqttClient.on('message', async (topic, message) => {
    const parts = topic.split('/');
    if (parts.length !== 3 || parts[2] !== 'gas') return;
    const houseCode = parts[1];

    let ppm;
    try {
      ppm = JSON.parse(message.toString()).ppm;
    } catch { return; }   // Ignorer les messages malformes

    let state = await HouseState.findOne({ houseCode });
    if (!state) state = await HouseState.create({ houseCode });

    const threshold = state.gas.threshold ?? 400;
    state.gas.level = ppm;
    await state.save();

    if (ppm >= threshold && !state.emergencyMode) {
      // ---- ACTIVATION DE L'URGENCE ----
      state.emergencyMode  = true;
      state.gas.lastAlert  = new Date();
      state.gas.valveOpen  = false;
      await state.save();

      // 1. Alerter tous les clients web de la maison
      io.to(houseCode).emit('house:gas-emergency', { ppm, threshold, houseCode });

      // 2. Fermer la vanne gaz
      mqttClient.publish(`command/${houseCode}/valve`, 'CLOSE', { qos: 1 });

      // 3. Activer le buzzer
      mqttClient.publish(`command/${houseCode}/alarm`, 'ON', { qos: 1 });

      // 4. Journaliser l'audit
      await AuditLog.create({
        actorName: 'Systeme',
        actorRole: 'system',
        action: `Urgence gaz : ${ppm.toFixed(0)} PPM depasse le seuil ${threshold}`,
        category: 'security',
        houseCode
      });
    }
  });
};
\end{lstlisting}

\section{Calcul PPM du Capteur Gaz MQ6}

Le capteur MQ6 produit une tension analogique proportionnelle à la concentration de gaz.
L'ESP32 lit cette valeur via l'ADC et la convertit en PPM en utilisant une formule
calibrée :

\begin{lstlisting}[style=codeblock,language=C++,caption={Calcul PPM Capteur Gaz (ESP32)}]
float readGasPPM() {
  int raw      = analogRead(GAS_PIN);           // 0 - 4095 (ADC 12 bits)
  float voltage= raw * 3.3f / 4095.f;
  float Rs     = ((3.3f - voltage) / voltage) * 10.0f;  // RL = 10 kOhm
  float ratio  = Rs / 9.83f;                            // calibration Ro
  return 1000.0f * pow(ratio, -2.37f);                  // PPM approximation
}

// Publier toutes les 3 secondes
if (millis() - lastGasPublish >= 3000) {
  float ppm = readGasPPM();
  String payload = "{\"ppm\":" + String(ppm, 1) + "}";
  mqttClient.publish(gasTopic.c_str(), payload.c_str());

  // Securite locale : declencher le buzzer si ADC brut > 800
  if (analogRead(GAS_PIN) > 800) tone(BUZZER_PIN, 1000);
  else                           noTone(BUZZER_PIN);
  lastGasPublish = millis();
}
\end{lstlisting}

\section{Service DHT11 et Firmware}

\begin{lstlisting}[style=codeblock,language=C++,caption={Lecture DHT11 et Publication MQTT (ESP32)}]
#include <DHT.h>
DHT dht(DHT_PIN, DHT11);

// Dans loop() :
if (millis() - lastDHTPublish >= 10000) {
  float temperature = dht.readTemperature();
  float humidity    = dht.readHumidity();

  if (!isnan(temperature) && !isnan(humidity)) {
    String payload = "{\"temperature\":" + String(temperature, 1)
                   + ",\"humidity\":"    + String(humidity,    1) + "}";
    mqttClient.publish(dhtTopic.c_str(), payload.c_str());
  }
  lastDHTPublish = millis();
}
\end{lstlisting}

\section{Overlay d'Urgence Frontend}

\begin{lstlisting}[style=codeblock,language=JavaScript,caption={Composant GasEmergencyOverlay}]
const GasEmergencyOverlay = ({ ppm, onClear, isAdmin }) => (
  <motion.div
    className="fixed inset-0 z-[9999] bg-red-600/90 flex flex-col
               items-center justify-center text-white"
    animate={{ opacity: [1, 0.8, 1] }}
    transition={{ repeat: Infinity, duration: 1.2 }}
  >
    <AlertTriangle size={96} className="mb-6 animate-bounce" />
    <h1 className="text-5xl font-bold mb-4">FUITE DE GAZ DETECTEE</h1>
    <p className="text-2xl mb-2">Niveau actuel : {ppm?.toFixed(0)} PPM</p>
    <p className="text-lg mb-8 opacity-80">
      Vanne fermee - Alarme active - Evacuez immediatement
    </p>
    {isAdmin && (
      <button onClick={onClear}
        className="px-8 py-3 bg-white text-red-700 rounded-full
                   font-bold text-lg hover:bg-red-100 transition">
        Lever l'Urgence
      </button>
    )}
  </motion.div>
);
\end{lstlisting}

\section{Revue et Rétrospective Sprint 4}

\begin{tcolorbox}[goalbox,title={Revue Sprint 4}]
\textbf{Complété :} 6 éléments sur 6 (36 PS). Urgence gaz vérifiée sur le matériel. \\
\textbf{Points forts :} Souffler sur le capteur MQ6 a déclenché l'overlay d'urgence en
500~ms, activé le buzzer physique, et le servo de vanne s'est fermé. L'admin a levé
l'urgence depuis le tableau de bord et le buzzer s'est tu en 1 seconde. \\
\textbf{Rétrospective :}
\begin{itemize}[leftmargin=1cm,itemsep=2pt]
  \item \textbf{Ce qui a bien fonctionné :} L'architecture de sécurité à deux couches
        (déclencheur local ESP32 ADC~$>$~800 + seuil serveur 400~PPM) fournit une
        défense en profondeur.
  \item \textbf{Défi :} Le capteur DHT11 retourne occasionnellement NaN lors des lectures
        à froid ; géré par un guard \texttt{isnan()} avant la publication MQTT.
  \item \textbf{Amélioration :} Ajout de \texttt{qos:1} aux commandes alarme et vanne
        pour survivre aux déconnexions MQTT brèves lors des scénarios d'urgence.
\end{itemize}
\end{tcolorbox}
